\section{Случайные события и вероятностное пространство}\subsection{Пространство элементарных исходов}Основной концепцией теории вероятностей является формализация случайного эксперимента.
\textbf{Пространство элементарных исходов} $\Omega$ — это множество всех возможных взаимно исключающих исходов $\omega$, которыми может завершиться эксперимент.Событие $A$ — это любое подмножество исходов $A \subseteq \Omega$.\subsection{Классическое определение вероятности}Если эксперимент удовлетворяет двум условиям:Число всех исходов $\vert{}\Omega\vert{} = n$ конечно.Все исходы равновозможны (симметрия эксперимента).То вероятность наступления события $A$, состоящего из $m$ исходов, определяется как отношение:
$$P(A) = \frac{\vert{}A\vert{}}{\vert{}\Omega\vert{}} = \frac{m}{n}.$$\subsection{Теоремы сложения вероятностей}\begin{theorem}[Теорема сложения для несовместных событий]
Если события $A$ и $B$ не могут наступить одновременно (т.е. их пересечение пусто: $A \cap B = \varnothing$), то вероятность того, что наступит хотя бы одно из них, равна сумме их вероятностей:
$$P(A \cup B) = P(A) + P(B).$$
\end{theorem}\begin{proof}
Пусть общее число исходов равно $n$. Событию $A$ благоприятствует $m_A$ исходов, событию $B$ — $m_B$ исходов. Поскольку $A$ и $B$ не пересекаются, ни один исход не благоприятствует обоим событиям одновременно. Значит, событию $A \cup B$ благоприятствует ровно $m_A + m_B$ исходов.
$$P(A \cup B) = \frac{m_A + m_B}{n} = \frac{m_A}{n} + \frac{m_B}{n} = P(A) + P(B).$$
\end{proof}\begin{theorem}[Общая теорема сложения вероятностей]
Для любых двух событий $A$ и $B$ (совместных или несовместных):
$$P(A \cup B) = P(A) + P(B) - P(A \cap B).$$
\end{theorem}\begin{proof}
Представим событие $A \cup B$ как объединение двух несовместных событий: самого события $A$ и той части события $B$, которая не входит в $A$. Формально: $A \cup B = A \cup (B \setminus A)$.
Так как эти события несовместны, применим предыдущую теорему:
$$P(A \cup B) = P(A) + P(B \setminus A). \quad (1)$$
В свою очередь, само событие $B$ можно разбить на две несовместные части: те исходы, которые входят только в $B$, и те, которые входят и в $B$, и в $A$.
$$B = (B \setminus A) \cup (A \cap B).$$
Вероятность $B$ равна сумме их вероятностей:
$$P(B) = P(B \setminus A) + P(A \cap B) \implies P(B \setminus A) = P(B) - P(A \cap B).$$
Подставим полученное выражение для $P(B \setminus A)$ в уравнение (1):
$$P(A \cup B) = P(A) + P(B) - P(A \cap B).$$
Теорема доказана.
\end{proof}\begin{theorem}[Формула включений-исключений для 3 событий]
$$P(A \cup B \cup C) = P(A) + P(B) + P(C) - P(A \cap B) - P(A \cap C) - P(B \cap C) + P(A \cap B \cap C).$$
\end{theorem}\begin{proof}
Обозначим $D = A \cup B$. Тогда $A \cup B \cup C = D \cup C$. Применим общую теорему сложения для двух событий $D$ и $C$:
$$P(D \cup C) = P(D) + P(C) - P(D \cap C).$$
Раскроем $P(D)$: $P(D) = P(A \cup B) = P(A) + P(B) - P(A \cap B)$. \
Теперь раскроем $P(D \cap C)$:
$$D \cap C = (A \cup B) \cap C = (A \cap C) \cup (B \cap C).$$
Вероятность этого объединения (по теореме сложения для двух событий):
$$P((A \cap C) \cup (B \cap C)) = P(A \cap C) + P(B \cap C) - P((A \cap C) \cap (B \cap C)).$$
Заметим, что $(A \cap C) \cap (B \cap C) = A \cap B \cap C$. \
Собираем все воедино:
\begin{align*}
P(A \cup B \cup C) &= [P(A) + P(B) - P(A \cap B)] + P(C) - [P(A \cap C) + P(B \cap C) - P(A \cap B \cap C)] \
&= P(A) + P(B) + P(C) - P(A \cap B) - P(A \cap C) - P(B \cap C) + P(A \cap B \cap C).
\end{align*}
\end{proof}